\chapter{Mastoiditis}
\index{Mastoiditis}

\section*{Definition and Overview}
Mastoiditis is a serious, potentially life-threatening bacterial infection of the mastoid air cells, which are the hollow, air-filled spaces located in the mastoid process of the temporal bone immediately behind the ear. It typically develops as a direct complication of untreated or inadequately treated acute otitis media (AOM), extending from the middle ear cavity into the surrounding bone. Mastoiditis can be acute or chronic, with the acute form characterised by rapid destruction of the bony septa separating the mastoid air cells (coalescent mastoiditis). Although the incidence of mastoiditis has declined significantly since the widespread introduction of antibiotics, it remains a critical clinical concern due to the risk of intracranial and extracranial complications.

\section*{Epidemiology}
Mastoiditis occurs predominantly in infants and young children under the age of two, whose Eustachian tubes are shorter, wider, and more horizontal, facilitating the spread of middle ear infections. However, it can affect individuals of any age. The incidence of acute mastoiditis is estimated at approximately 1.5 to 4 cases per 100,000 children per year in developed countries like some countries. The prevalence is slightly higher in males and during winter months, mirroring the seasonal trends of acute otitis media. Immunocompromised individuals and those with chronic middle ear pathology are at increased risk.

\section*{Aetiology and Pathophysiology}
The development of mastoiditis involves the progression of a middle ear infection into the adjacent temporal bone.
\begin{itemize}
    \item \textbf{Middle ear extension:} Bacteria from an acute otitis media episode migrate through the aditus ad antrum (the narrow channel connecting the middle ear cavity to the mastoid antrum).
    \item \textbf{Inflammatory blockage:} Swelling of the mucosal lining blocks the aditus ad antrum, trapping purulent exudate within the mastoid air cells and preventing normal Eustachian tube drainage.
    \item \textbf{Pressure necrosis:} Accumulation of pus leads to increased pressure, causing localised ischaemia, mucosal necrosis, and the dissolution of the thin bony trabeculae separating the air cells, a process known as coalescent mastoiditis.
    \item \textbf{Causative pathogens:} The most common pathogen isolated is Streptococcus pneumoniae, followed by Streptococcus pyogenes, Staphylococcus aureus, and Pseudomonas aeruginosa (especially in chronic cases).
    \item \textbf{Osteomyelitis:} In severe cases, the infection spreads from the air cells to the cortical bone of the temporal bone, leading to subperiosteal abscess formation or intracranial extension.
\end{itemize}

\section*{Clinical Features}
Mastoiditis presents with distinct localised signs and symptoms, usually appearing several days to weeks after the onset of acute otitis media.
\begin{itemize}
    \item \textbf{Post-auricular signs:} Redness, swelling, warmth, and marked tenderness over the mastoid process behind the ear.
    \item \textbf{Protruding ear:} Displacement of the auricle (pinna) outward and downward due to swelling or subperiosteal abscess formation.
    \item \textbf{Otorrhoea and otalgia:} Deep, throbbing ear pain (otalgia) and purulent ear discharge (otorrhoea), often following tympanic membrane rupture.
    \item \textbf{Systemic toxicity:} High fever, irritability, lethargy, and poor feeding, especially in infants.
    \item \textbf{External auditory canal changes:} Sagging of the posterosuperior wall of the external auditory canal due to adjacent mastoid inflammation.
\end{itemize}

\section*{Differential Diagnosis}
It is critical to distinguish mastoiditis from other conditions presenting with ear pain and peri-auricular swelling.
\begin{itemize}
    \item \textbf{Post-auricular lymphadenitis:} Inflammation and enlargement of the lymph nodes behind the ear, usually secondary to a scalp or skin infection, without middle ear disease or ear displacement.
    \item \textbf{Severe otitis externa:} Infection of the external ear canal that can cause severe pain, swelling, and redness of the pinna, but without middle ear effusion or bony mastoid tenderness.
    \item \textbf{Mumps (infectious parotitis):} Viral infection causing painful swelling of the parotid salivary glands, located anterior and inferior to the ear, rather than posterior.
    \item \textbf{Histiocytosis or neoplasm:} Rare non-infectious lesions of the temporal bone (such as Langerhans cell histiocytosis or rhabdomyosarcoma) that can mimic mastoiditis on imaging.
    \item \textbf{Cholesteatoma:} A benign keratinising squamous epithelial growth in the middle ear that can erode bone and present with chronic, foul-smelling drainage and mastoiditis-like complications.
\end{itemize}

\section*{When to Seek Medical Care}
Any child or adult showing signs of ear pain accompanied by swelling or redness behind the ear requires urgent medical evaluation by a general practitioner or emergency department.

\textbf{Contact your local emergency services for:}
\begin{itemize}
    \item Signs of intracranial spread, including altered mental state, stiff neck (meningismus), severe headache, persistent vomiting, or seizures.
    \item Focal neurological deficits, such as facial nerve paralysis (inability to move muscles on one side of the face) or double vision.
    \item Extreme lethargy, floppy baby syndrome, or difficulty waking up.
\end{itemize}

\section*{Diagnosis and Investigations}
Diagnosis is based on clinical presentation, supported by laboratory tests, cultures, and imaging studies.
\begin{itemize}
    \item \textbf{Otoscopy:} Visual inspection of the tympanic membrane, which is typically bulging, erythematous, and immobile, or shows a perforation with active purulent discharge.
    \item \textbf{CT scan of the temporal bones:} The imaging modality of choice, revealing opacification of the mastoid air cells and loss of the bony septa (coalescence), and ruling out intracranial complications.
    \item \textbf{Ear swab and culture:} Collection of middle ear fluid via tympanocentesis (if the eardrum is intact) or from active drainage to identify the specific bacterial pathogen and determine antibiotic sensitivity.
    \item \textbf{Inflammatory markers:} Full blood count showing leukocytosis, and elevated C-reactive protein (CRP) or erythrocyte sedimentation rate (ESR).
    \item \textbf{MRI scan:} Indicated if intracranial complications (such as sigmoid sinus thrombosis, brain abscess, or meningitis) are suspected, providing superior soft-tissue detail.
\end{itemize}

\section*{Management}
Management requires prompt hospitalisation, intravenous antibiotics, and often surgical intervention to drain the infection.
\begin{itemize}
    \item \textbf{Intravenous antibiotics:} Empiric broad-spectrum intravenous therapy (such as ceftriaxone or cefotaxime) is initiated immediately and tailored once culture results are available.
    \item \textbf{Myringotomy and tympanostomy tube insertion:} A minor surgical procedure to make a small incision in the eardrum to drain fluid, relieve pressure, and allow instillation of antibiotic drops.
    \item \textbf{Mastoidectomy:} Surgical removal of the infected mastoid air cells and cortical bone (simple mastoidectomy) is performed if there is a subperiosteal abscess, coalescent mastoiditis, or failure to improve after 24 to 48 hours of IV antibiotics.
    \item \textbf{Pain and fever management:} Administration of paracetamol and ibuprofen to manage pain and control high fevers.
    \item \textbf{Fluid replacement:} Intravenous hydration for patients who are unable to tolerate oral intake due to severe lethargy or vomiting.
\end{itemize}

\section*{Complications}
Mastoiditis is associated with severe complications if the infection spreads beyond the mastoid air cells.
\begin{itemize}
    \item \textbf{Subperiosteal abscess:} Accumulation of pus beneath the periosteum of the temporal bone, causing marked swelling and displacement of the ear.
    \item \textbf{Facial nerve palsy:} Damage or compression of the facial nerve (cranial nerve VII) as it traverses the temporal bone, leading to facial weakness or paralysis.
    \item \textbf{Labyrinthitis:} Extension of the infection into the inner ear, causing severe vertigo, sensorineural hearing loss, and tinnitus.
    \item \textbf{Meningitis and brain abscess:} Spreading of the infection into the cranial cavity, leading to life-threatening inflammation of the meninges or collection of pus in the brain.
    \item \textbf{Dural venous sinus thrombosis:} Septic thrombosis of the sigmoid or transverse dural venous sinuses, causing high intracranial pressure and neurological deficits.
    \item \textbf{Bezold's abscess:} A rare complication where pus tracks downward into the sternocleidomastoid muscle sheath in the neck.
\end{itemize}

\section*{Prognosis and Outcomes}
With early diagnosis and aggressive antibiotic treatment, the prognosis for acute mastoiditis is generally excellent. Most patients recover completely without long-term hearing loss or neurological sequelae. If surgical drainage is required, healing of the surgical site is usually complete within several weeks. However, the prognosis is guarded if serious intracranial complications have occurred, requiring prolonged neurosurgical and infectious disease management.

\section*{Prevention and Self-Care}
Preventative efforts focus on prompt treatment of ear infections and maintaining up-to-date immunisations.
\begin{itemize}
    \item \textbf{Prompt treatment of otitis media:} Seek medical attention for children with ear pain, fever, or fluid draining from the ear to ensure appropriate treatment.
    \item \textbf{Pneumococcal vaccination:} Ensure children receive the routine pneumococcal conjugate vaccine (PCV), which has significantly reduced the incidence of AOM and mastoiditis.
    \item \textbf{Influenza vaccination:} Annual flu vaccines reduce the risk of viral upper respiratory tract infections that can lead to secondary ear infections.
    \item \textbf{Avoid second-hand smoke:} Exposure to cigarette smoke increases the frequency and severity of middle ear infections in children.
    \item \textbf{Promote breastfeeding:} Breastfeeding for the first six months of life provides antibodies that protect infants against middle ear infections.
\end{itemize}

\section*{Sources and Further Reading}
The following organisations provide reliable information about mastoiditis and ear infections:
\begin{itemize}
    \item Royal Children's Hospital (Melbourne). \textit{Clinical guidelines on acute otitis media and its complications.} https://www.rch.org.au/
    \item Healthdirect Australia. \textit{Consumer guide to mastoiditis, symptoms, diagnosis, and treatment.} https://www.healthdirect.gov.au/mastitis-ear-infection
    \item NHS. \textit{Overview of mastoiditis, including causes, symptoms, and surgical procedures.} https://www.nhs.uk/conditions/mastoiditis/
    \item Mayo Clinic. \textit{Information on otitis media and potential complications like mastoiditis.} https://www.mayoclinic.org/diseases-conditions/ear-infection/symptoms-causes/syc-20351616
    \item Stanford Medicine (Ear, Nose \& Throat). \textit{Clinical information on the management of acute and chronic mastoiditis.} https://med.stanford.edu/ohns/
\end{itemize}
